\documentclass{article}

\usepackage{../packages}
\usepackage{../macros}

\title{$\kappa$-geometric classifying topoi}
\author{mark}

\begin{document}
\maketitle

In the course of understanding light condensed sets, and similar flavour of topoi, it is worthwhile to consider large sites, which nonetheless are essentially the classifying topos
of an algebraic theory - in some sense.

For instance we can consider the classifying topos of boolean algebras, given by presheaves $\Set[\mathrm{Bool}] := [\mathrm{Bool}_{\mathrm{fp}}, \Set]$.
This has a universal boolean algebra $U : \mathrm{Bool}_{\mathrm{fp}} \to \Set$, given by the forgetful functor.
This classifies boolean algebras in the sense that it is the image of the identity under the natural isomorphism:
\[ \mathrm{Geom}(\Set[\mathrm{Bool}], \Set[\mathrm{Bool}]) \simeq \mathrm{Bool}(\Set[\mathrm{Bool}]) \] 

Consider the category of light condensed sets, which is equivalently presented as presheaves $[\mathrm{Bool}_{\mathrm{cp}}, \Set]$.
This isn't a classifying topos of anything reasonable; in an answer on stack exchange, Peter Scholze finds a theory 
this topos classifies, but it isn't particularly interesting.

Instead we should imagine that this topos also really classifies the theory of boolean algebras; however the notion of morphism between topoi is not 
strong enough to make this precise, as they don't preserve the structure of countably presented boolean algebras.

Instead we extend to $\kappa$-geometric morphisms, which allow this correspondence to be made precise.
Let $\kappa$ be a cardinality/arity class. 

\begin{definition}
  Then a $\kappa$-geometric morphism $f : \mathscr{E} \to \mathscr{D}$ is a functor 
  $f^\ast : \mathscr{D} \to \mathscr{E}$ preserving all small colimits, and $\kappa$-ary limits.
\end{definition}

We will consider the category $\mathbf{Top}$ of Grothendieck topoi, with $\kappa$-geometric morphisms.

\begin{lemma}
  $\text{TODO}$ is the terminal $\kappa$-complete topos.
\end{lemma}
\begin{proof}
  For any $\mathscr{E}$ we get a functor $\Set \to \mathscr{E}$ given by the usual functor $X \mapsto \bigsqcup_{x \in X} 1$
  But this doesn't preserve colimits.. 
  So $\Set$ definitely isn't the terminal topos. 
\end{proof}


\section{Classifying morphisms into sheaf topoi}

\begin{definition}
  A functor $A : \cC \to \Set$ is called $\kappa$-flat if the left Kan extension along the Yoneda embedding $- \otimes_\cC A : \Psh{\cC} \to \Set$ preserves $\kappa$-ary limits.
\end{definition}

By definition $\kappa$-flat functors $A : \cC \to \Set$ correspond precisely to $\kappa$-geometric morphisms $\Set \to \Psh{\cC}$.

\begin{definition}
  A category $I$ is called $\kappa$-filtering if all $\kappa$-ary diagrams in $I$ have a cone.

  A functor $A : \cC \to \Set$ is called $\kappa$-filtering if the category of elements $\ast \downarrow A$ is $\kappa$-filtering.
\end{definition}

\begin{lemma}
  A diagram $D : I \to \ast \downarrow A$ is precisely determined by:
  \begin{itemize}
    \item A collection $c_i \in \cC$ for $i \in I$. 
    \item An element $a \in \lim_I A(c_i)$, where we consider the projection diagram in $\cC$.
  \end{itemize} 
\end{lemma}

\begin{proposition}
  A functor $A : \cC \to \Set$ is $\kappa$-flat iff it is $\kappa$-filtering.
\end{proposition}
\begin{proof}
  ($\Rightarrow$) Suppose $A$ is $\kappa$-flat. 
  Let $D : I \to \ast \downarrow A$ be a $\kappa$-ary diagram.
  Let $c_i : \cC$ be the projection fo $D(i)$ for $i \in I$.
  Let $a_i : A(c_i)$ be the second projection of $D(i)$.
  Then as $A$ preserves $\kappa$-ary limits we have an isomorphism
  \[ \left ( \lim_{i \in I} \yo c_i \right ) \otimes A \to \lim_{i \in I} (\yo c_i \otimes A) \simeq \lim_{i \in I} A(c_i) \] 
  The map is given by 
  \[  w \otimes b \mapsto w_i \cdot b \]
  Since by assumption for any $f : i \to j$ we have $f_\ast(a_i) = a_j$ we can see that the collection $(a_i)$ forms an element of $\lim_{i \in I} A(c_i)$. 
  Hence there is $b \in \cC$ and $w_i : b \to c_i$ so that $w_i \cdot b = a_i$ for each $i \in I$.
  This is precisely a cone over $D$.
  Hence $A$ is $\kappa$-filtering.

  ($\Leftarrow$) Suppoes $A$ is $\kappa$-filtering. Note by standard arguments in SIGAL, $A$ must be flat. 
  Hence to show it is $\kappa$ flat it suffices to show that $- \otimes A$ preserves $\kappa$-ary products. \todo{Can't figure this out constructively. Would be easy if we new $\kappa$-filtered colimits and $\kappa$-finite limits commuted. This is relatively straightforward with $\kappa$-ary choice. Not sure if there's a constructive version.}
\end{proof}


\begin{proposition}
  Suppose $\cC$ is $\kappa$-complete. Then a functor $A : \cC \to \Set$ is $\kappa$-flat iff $A$ is $\kappa$-exact.
\end{proposition}
\begin{proof}
  ($\Rightarrow$) Suppose $A$ is $\kappa$-flat. Then $- \otimes_\cC A$ preserves all $\kappa$-ary limits.
  But then, as the yoneda embedding preserves all limits we have 
  \[ A \simeq (- \otimes_\cC A) \circ \yo : \cC \to \Set \]
  preserves all $\kappa$-ary limits in $\cC$. 
  Hence $A$ is $\kappa$-complete.

  ($\Leftarrow$) 
  Suppose $A$ is $\kappa$-complete. 
  Now suppose $D : I \to \int_{\cC} A$ is a $\kappa$-ary diagram. 
  Then let $c_i$ be $\pi D(i)$ and $a_i$ be the second projection of $D(i)$. 
  Then we construct a cone over the diagram to be the object with first projection $\lim_I c_i$ which exists as $\cC$ is $\kappa$-complete.
  Since $A$ preserves $\kappa$-ary limits we have 
  $A(\lim_I c_i) \simeq \lim_I A(c_i)$.
  The collection $a_i \in A(c_i)$ satisfy the condition to form an element of the limit by construction. 
  Hence we have a single element $a \in A(\lim_I c_i)$ satisfying $\pi_i \cdot a = a_i$. 
  Thus this is a cone in the diagram category.
  Hence $A$ is $\kappa$-filtering.
  Thus by the above $A$ is $\kappa$-flat.
\end{proof}




\section{Doctrines}

\begin{definition}
  A doctrine $\mathbb{D}$ is an essentially small class of categories. 
\end{definition}



\end{document}